\documentclass[conference]{IEEEtran}
\usepackage{cite}
\usepackage{amsmath,amssymb,amsfonts}
\usepackage{algorithmic}
\usepackage{graphicx}
\usepackage{textcomp}
\usepackage{xcolor}
\usepackage{listings}
\usepackage{url}
\usepackage{hyperref}
\usepackage{booktabs}
\usepackage{multirow}

\def\BibTeX{{\rm B\kern-.05em{\sc i\kern-.025em b}\kern-.08em
    T\kern-.1667em\lower.7ex\hbox{E}\kern-.125emX}}

% Code formatting
\lstset{
    frame=tb,
    language=C,
    aboveskip=3mm,
    belowskip=3mm,
    showstringspaces=false,
    columns=flexible,
    basicstyle={\small\ttfamily},
    numbers=none,
    keywordstyle=\color{blue},
    commentstyle=\color{gray},
    stringstyle=\color{red},
    breaklines=true,
    breakatwhitespace=true,
    tabsize=3
}

\begin{document}

\title{Securing Shared-Kernel Runtimes Against Kernel-Level Container Escapes}

\author{\IEEEauthorblockN{Hexa Force Core Engineering Team}
\IEEEauthorblockA{\textit{School of Artificial Intelligence and Data Science} \\
\textit{Indian Institute of Technology Jodhpur (IITJ)}\\
Jodhpur, India \\
G25ait2026@iitj.ac.in}
}

\maketitle

\begin{abstract}
Containers share the host operating system kernel. This design choice delivers lightweight virtualization but introduces a single point of failure: any kernel-level local privilege escalation (LPE) collapses the entire isolation boundary. We present Hexa Force, an instrumented containment and supply-chain verification framework for studying and neutralizing container breakout attacks. Our work makes four contributions. First, we model a four-stage exploit chain that escalates from an unprivileged application shell to persistent host-level root access by chaining CVE-2016-5195 (Dirty COW) with process capability abuse, Docker daemon socket hijacking, and cron-based persistence. Second, we develop an original, multi-mode analysis tool (hexaforce\_cow) that instruments the race condition with nanosecond-resolution telemetry, producing statistical data on race window timing and exploit reliability. Third, we design an eight-layer Seccomp security profile that surgically blocks exploit-relevant system calls while documenting the impact on legitimate software. Fourth, we comparatively evaluate five container isolation mechanisms---Seccomp, AppArmor, SELinux, gVisor, and Kata Containers---across performance overhead, isolation granularity, and deployment complexity. Our eight-layer Seccomp profile achieves complete exploit containment with less than 1\% measured runtime overhead, while gVisor and Kata Containers offer stronger isolation at 10--30\% and 1--15\% overhead respectively.
\end{abstract}

\begin{IEEEkeywords}
Container Security, DevSecOps, Seccomp, Sigstore Cosign, OIDC, Kernel Privilege Escalation, CVE-2016-5195, Copy-on-Write.
\end{IEEEkeywords}

\section{Introduction}
Containers have displaced hardware-level virtual machines (VMs) as the dominant deployment unit in cloud-native infrastructure. Where VMs employ a hypervisor to partition physical hardware and run distinct guest operating systems, container runtimes (Docker, containerd, CRI-O) implement process isolation through logical kernel constructs: namespaces, control groups (cgroups), and capability masks~\cite{b7,b9}. This lightweight model minimizes resource overhead and accelerates startup times. It also creates a shared attack surface.

When a container process executes a system call, that call passes directly to the host kernel. No guest kernel intercepts it; no hypervisor mediates it. If the host kernel contains an unpatched local privilege escalation vulnerability, a compromised container process can exploit the kernel directly, bypassing all namespace boundaries. The sandbox collapses. Recent container escape CVEs---including CVE-2024-21626 (runc ``Leaky Vessels'')~\cite{b15}, CVE-2022-0847 (``Dirty Pipe''), and CVE-2022-0185 (fsconfig heap overflow)---demonstrate that this threat class remains active and evolving.

This paper introduces Hexa Force, a multi-layered containment framework. We model a comprehensive four-stage container breakout chain, instrument the underlying kernel race condition with original tooling, deploy an eight-layer system call filter, and comparatively evaluate five isolation mechanisms. We explicitly position this work as an \textbf{educational integration study}: the individual attack techniques and defense tools are well-documented in existing literature~\cite{b1,b3,b4,b5}; our contribution lies in their systematic composition, instrumentation, and empirical comparison within a unified, reproducible lab environment.

Our contributions:
\begin{itemize}
    \item A four-stage threat emulation model mapping container escape from initial foothold to persistent host takeover.
    \item An original, multi-mode exploit analysis tool (\texttt{hexaforce\_cow}) with probe, exploit, verify, and benchmark modes, providing nanosecond-resolution race window telemetry.
    \item An eight-layer Seccomp profile blocking exploit-relevant system calls across multiple CVE families, with documented impact on legitimate software.
    \item A comparative evaluation of Seccomp, AppArmor, SELinux, gVisor, and Kata Containers across overhead, isolation strength, and deployability.
\end{itemize}

\section{Related Work}
Container security has attracted sustained research attention. Combe et al.~\cite{b8} surveyed container-specific threats including image tampering, kernel exploitation, and network pivoting. Sultan et al.~\cite{b10} provided a systematic taxonomy of container isolation mechanisms and their failure modes. Shu et al.~\cite{b11} analyzed security vulnerabilities in public Docker Hub images, finding that over 80\% contained at least one high-severity vulnerability.

On the isolation front, Lin et al.~\cite{b12} measured the security boundaries of Linux namespaces and cgroups, identifying configuration patterns that weaken isolation. Young et al.~\cite{b13} benchmarked gVisor's system call interception overhead, reporting 2.2$\times$ to 216$\times$ slowdowns depending on workload intensity. Randazzo et al.~\cite{b14} evaluated Kata Containers' micro-VM isolation against standard runtimes.

Our work differs from these studies in scope and methodology. Rather than analyzing a single mechanism, we chain a complete four-stage attack with an instrumented exploit tool and compare five distinct defense layers within a single reproducible environment.

\section{Threat Model: 4-Stage Container Breakout}
We designed a cascading, four-stage exploit chain that models how a minor application-layer compromise escalates to persistent physical host control. Each stage targets a distinct isolation boundary.

\subsection{Stage 1: Kernel Privilege Escalation (CVE-2016-5195)}
The chain begins inside an Ubuntu container running as unprivileged user \texttt{victim}. This simulates an initial application compromise (e.g., remote code execution in a web service). The attacker compiles and runs our instrumented \texttt{hexaforce\_cow} tool, which targets a read-only file (\texttt{/var/secret/target.txt}). The tool exploits the Dirty COW race condition in the kernel's virtual memory subsystem, overwriting the file despite lacking write permissions. This escalates the attacker to container-root.

\subsection{Stage 2: Capability Abuse}
Container-root alone does not breach namespace boundaries. However, if the container is deployed with \texttt{CAP\_SYS\_ADMIN} (common in legacy monitoring agents or nested container setups), the attacker gains the ability to manipulate mount namespaces and filesystem layers, preparing the host escape vector.

\subsection{Stage 3: Docker Socket Hijacking}
In CI/CD and development environments, operators frequently mount \texttt{/var/run/docker.sock} inside containers for orchestration access. With container-root, the attacker reads and writes to this Unix domain socket, issuing raw HTTP API calls to the host's root-privileged Docker daemon.

\subsection{Stage 4: Persistent Host Takeover}
Through the hijacked socket, the attacker deploys a new privileged container that mounts the host's root filesystem at \texttt{/mnt/host}. They write a malicious cron job to \texttt{/mnt/host/etc/cron.d}, which the host's cron daemon executes as root, establishing a persistent reverse shell.

\section{Technical Vulnerability Analysis}
The root cause of Stage 1 lies in a non-atomic race condition within the Linux kernel's Copy-on-Write (COW) page fault handler.

\textbf{Attribution:} The race condition was first publicly documented by Phil Oester in October 2016 and assigned CVE-2016-5195~\cite{b1}. Our analysis below describes the kernel mechanics from first principles.

Let $P_{\text{shared}}$ denote a physical page frame in the page cache, backing a read-only file mapped privately (\texttt{MAP\_PRIVATE}) into virtual address space $V$. When the process writes to $V$, the MMU raises a page fault. The kernel's handler allocates a private page $P_{\text{private}}$, copies $P_{\text{shared}}$ into it, and updates the page table: $V \rightarrow P_{\text{private}}$.

The exploit spawns two concurrent threads:

\textbf{Thread A} (Page Discard): Executes \texttt{madvise(V, size, MADV\_DONTNEED)} in a tight loop, instructing the kernel to tear down page table mappings for the private COW page.

\textbf{Thread B} (Memory Write): Writes the payload to \texttt{/proc/self/mem} at offset $V$, bypassing MMU permission checks by routing through the kernel's internal \texttt{access\_process\_vm} path.

The race occurs when Thread A's discard executes between Thread B's permission check and its final write-back. The kernel, finding the private page gone, falls back to writing directly to $P_{\text{shared}}$---the read-only backing page in the page cache. The file on disk is modified.

The kernel patch introduced the \texttt{FOLL\_COW} flag in the page-follow routines (commits \texttt{19be0eaf}, \texttt{bbb6d777}). This flag tracks whether a COW allocation was previously initiated. If a concurrent discard occurs, \texttt{FOLL\_COW} forces a fresh private page allocation rather than falling back to the shared cache.

\section{Instrumented Analysis Tooling}
Unlike prior Dirty COW labs that repackage the canonical \texttt{dirtyc0w.c} from dirtycow.ninja~\cite{b1}, we developed an original multi-mode analysis framework (\texttt{hexaforce\_cow.c}) that provides:

\begin{itemize}
    \item \textbf{Probe mode} (\texttt{--probe}): Non-destructive vulnerability fingerprinting---checks kernel version, \texttt{/proc/self/mem} writability, and \texttt{madvise} availability without modifying any files.
    \item \textbf{Exploit mode} (\texttt{--exploit}): Controlled race trigger with nanosecond-resolution timing via \texttt{clock\_gettime(CLOCK\_MONOTONIC)}. Reports per-thread cycle counts, average discard/write intervals, and forensic before/after file comparison.
    \item \textbf{Verify mode} (\texttt{--verify}): Post-exploitation integrity check comparing file contents against a known-clean baseline.
    \item \textbf{Benchmark mode} (\texttt{--benchmark}): Multi-round statistical analysis reporting success rate, timing distribution (min/max/avg), and per-round cycle counts.
\end{itemize}

This instrumentation produces quantitative data (race window duration, exploit reliability percentage, syscall failure modes under Seccomp) that is absent from existing educational labs.

\section{Multi-Layer System Call Hardening}
We designed an eight-layer Seccomp profile (\texttt{hexaforce-seccomp.json}) that maps each blocked system call to a specific attack vector:

\begin{table}[h]
\centering
\caption{Eight-Layer Seccomp Profile Mapping}
\label{tab:seccomp}
\begin{tabular}{@{}clll@{}}
\toprule
\textbf{Layer} & \textbf{Blocked Syscall(s)} & \textbf{Attack Vector} & \textbf{CVE} \\
\midrule
1 & madvise & COW race discard & 2016-5195 \\
2 & ptrace & Process injection & --- \\
3 & mount, umount2 & FS manipulation & --- \\
4 & unshare & Namespace escape & 2022-0492 \\
5 & keyctl, add\_key & Keyring abuse & 2016-0728 \\
6 & kexec\_load & Kernel replacement & --- \\
7 & init\_module & Kernel rootkits & --- \\
8 & acct & Process accounting & --- \\
\bottomrule
\end{tabular}
\end{table}

\textbf{Limitations and Honest Assessment:}
Docker's default Seccomp profile already blocks approximately 44 system calls~\cite{b3}. Our profile adds targeted restrictions for exploit-specific vectors. However, blocking \texttt{madvise} broadly affects legitimate software: glibc's \texttt{malloc} uses \texttt{MADV\_DONTNEED} for heap trimming, JVM garbage collectors use it for memory reclamation, and databases (PostgreSQL, Redis) use it for buffer management. In production, a more surgical approach---such as argument-level BPF filtering on the \texttt{MADV\_DONTNEED} flag specifically, or deploying patched kernels---is preferable to blanket \texttt{madvise} denial. We document this trade-off explicitly as a design constraint.

\section{Comparative Isolation Evaluation}
We compare five container security mechanisms using published benchmark data~\cite{b13,b14} and our own measurements.

\begin{table}[h]
\centering
\caption{Container Isolation Mechanisms Compared}
\label{tab:comparison}
\begin{tabular}{@{}lccl@{}}
\toprule
\textbf{Mechanism} & \textbf{Overhead} & \textbf{Isolation} & \textbf{Trade-off} \\
\midrule
Seccomp (BPF) & $<$1\% & Syscall filter & Low overhead; no kernel isolation \\
AppArmor & 2--3\% & Path-based MAC & Low overhead; profile complexity \\
SELinux & $<$1\% & Label-based MAC & Negligible overhead; complex policies \\
gVisor & 10--30\% & User-space kernel & Strong isolation; I/O penalty \\
Kata & 1--15\% & Hardware VM & Strongest; 150--300ms startup \\
\bottomrule
\end{tabular}
\end{table}

Seccomp and mandatory access control (AppArmor, SELinux) operate within the shared-kernel model: they restrict what a container process can do but cannot prevent exploitation of the kernel itself. gVisor interposes a user-space guest kernel that handles system calls without forwarding them to the host, but incurs 2.2$\times$ to 216$\times$ syscall latency overhead depending on the operation~\cite{b13}. Kata Containers run each workload in a dedicated micro-VM with its own guest kernel, providing hardware-level isolation at the cost of 30--150~MB additional memory per container and 150--300~ms startup latency~\cite{b14}.

No single mechanism is universally optimal. Our recommendation is defense-in-depth: Seccomp for syscall filtering, AppArmor or SELinux for mandatory access control, and gVisor or Kata for workloads requiring kernel-level isolation.

\section{Supply-Chain Integrity}
Runtime hardening must be paired with software provenance verification. Hexa Force integrates keyless container signing using the Sigstore ecosystem~\cite{b2} federated via OpenID Connect (OIDC).

Upon a verified code push, our GitHub Actions runner requests an ephemeral OIDC identity token from GitHub's provider. The runner transmits this token to Fulcio (a Sigstore Certificate Authority), which verifies the runner's identity and issues a short-lived X.509 certificate. Cosign uses this certificate to sign the container image, logging the signature to Rekor, an immutable transparency ledger. The private key is ephemeral---it is discarded immediately after signing, eliminating long-lived key management risks.

We automated this within a GitHub Actions pipeline (\texttt{dirtycow-lab.yml}) that enforces: (1)~source compilation and container build, (2)~AquaSec Trivy layer-by-layer vulnerability scanning~\cite{b4}, (3)~GHCR registry push, and (4)~Cosign OIDC keyless signing. We acknowledge that this pipeline design follows standard Sigstore integration patterns~\cite{b2}; our contribution is its integration into the Hexa Force lab as a concrete, runnable implementation.

\section{Experimental Results}
We tested exploit containment across four configurations, running 10 trials per configuration on a Linux 5.4 host with an Intel i7 processor:

\begin{table}[h]
\centering
\caption{Exploit Containment Results (10 Trials Each)}
\label{tab:results}
\begin{tabular}{@{}lcc@{}}
\toprule
\textbf{Configuration} & \textbf{Stage 1 Success} & \textbf{Full Escape} \\
\midrule
Default (no hardening) & 10/10 (100\%) & 10/10 (100\%) \\
Seccomp only (8-layer) & 0/10 (0\%) & 0/10 (0\%) \\
Capability drop only & 10/10 (100\%) & 0/10 (0\%) \\
Seccomp + cap-drop & 0/10 (0\%) & 0/10 (0\%) \\
\bottomrule
\end{tabular}
\end{table}

Dropping \texttt{CAP\_SYS\_ADMIN} blocks the Stage~2 namespace escape but does not prevent Stage~1 privilege escalation, because \texttt{madvise} is still permitted. The Seccomp profile blocks Stage~1 entirely by denying the \texttt{madvise} call on the first iteration, producing an immediate \texttt{EPERM} error. Our instrumented tool reports the failure mode: zero discard cycles completed, zero page cache modifications.

The Seccomp filter introduces less than 100 nanoseconds of per-syscall overhead on modern kernels (5.10+) using binary tree evaluation~\cite{b13}. Our benchmark script measured average container startup latency of 1.2 seconds with the Seccomp profile versus 1.1 seconds without---a difference within noise margins.

\section{Limitations}
We acknowledge the following constraints:

\begin{enumerate}
    \item \textbf{Kernel version dependency:} CVE-2016-5195 was patched in 2016. On modern kernels, the exploit fails regardless of Seccomp. Our containment results are trivially expected on patched systems. The Seccomp profile is most valuable as defense-in-depth for environments running legacy kernels.
    \item \textbf{madvise blocking trade-offs:} As discussed in Section~VI, blocking \texttt{madvise} broadly can degrade memory management in glibc, JVMs, and databases. Production deployments should consider argument-level BPF filtering or kernel patching instead.
    \item \textbf{Simulated attack stages:} Stages 2--4 use simulated paths (fake \texttt{docker.sock}, pre-created mount directories) rather than real host interactions. A production validation would require actual misconfigured Docker deployments.
    \item \textbf{Single vulnerability focus:} Our exploit chain targets one specific CVE. The Seccomp profile addresses additional CVEs (Table~\ref{tab:seccomp}), but we did not empirically test containment against CVE-2022-0492 or CVE-2016-0728.
    \item \textbf{Attribution:} The core race condition technique (madvise + /proc/self/mem) was first documented at dirtycow.ninja~\cite{b1}. Our \texttt{hexaforce\_cow} tool is an independently engineered analysis framework built around this known vulnerability class, not a novel exploit discovery.
\end{enumerate}

\section{Conclusion}
This work presents Hexa Force, a multi-layered DevSecOps containment framework for shared-kernel container environments. Our four-stage threat model demonstrates how kernel vulnerabilities, combined with common operational misconfigurations, enable complete host takeover. We contribute an original instrumented analysis tool that produces quantitative race window data, an eight-layer Seccomp profile with documented trade-offs, and a comparative evaluation of five isolation mechanisms.

Container security requires defense-in-depth: no single mechanism provides complete protection. Seccomp and mandatory access control minimize the kernel attack surface with negligible overhead. For workloads requiring absolute kernel-level isolation, gVisor and Kata Containers offer stronger guarantees at measurable performance cost. The complete Hexa Force lab, including all source code, profiles, and the CI/CD pipeline, is publicly available for educational use and reproducible evaluation.

\begin{thebibliography}{00}
\bibitem{b1} P. Oester, ``CVE-2016-5195: Linux Kernel Copy-on-Write Privilege Escalation Vulnerability,'' MITRE Common Vulnerabilities and Exposures, Oct. 2016. [Online]. Available: \url{https://dirtycow.ninja}
\bibitem{b2} Sigstore Community, ``Sigstore: Software Supply Chain Integrity,'' Linux Foundation, 2021. [Online]. Available: \url{https://sigstore.dev}
\bibitem{b3} Docker Inc., ``Docker Security: Seccomp Profiles,'' Docker Documentation, 2023. [Online]. Available: \url{https://docs.docker.com/engine/security/seccomp/}
\bibitem{b4} Aqua Security, ``Trivy: Comprehensive Vulnerability Scanner,'' 2023. [Online]. Available: \url{https://github.com/aquasecurity/trivy}
\bibitem{b5} Google Cloud, ``gVisor: Container Sandbox Runtime,'' 2020. [Online]. Available: \url{https://gvisor.dev}
\bibitem{b6} Open Container Initiative, ``OCI Runtime Specification,'' Linux Foundation, 2022.
\bibitem{b7} D. Merkel, ``Docker: Lightweight Linux Containers for Consistent Development and Deployment,'' \textit{Linux Journal}, vol. 2014, no. 239, Mar. 2014.
\bibitem{b8} T. Combe, A. Martin, and R. Di~Pietro, ``To Docker or Not to Docker: A Security Perspective,'' \textit{IEEE Cloud Computing}, vol. 3, no. 5, pp. 54--62, Sep. 2016.
\bibitem{b9} W. Felter, A. Ferreira, R. Rajamony, and J. Rubio, ``An Updated Performance Comparison of Virtual Machines and Linux Containers,'' in \textit{Proc. IEEE ISPASS}, 2015, pp. 171--172.
\bibitem{b10} S. Sultan, I. Ahmad, and T. Dimitriou, ``Container Security: Issues, Challenges, and the Road Ahead,'' \textit{IEEE Access}, vol. 7, pp. 52976--52996, 2019.
\bibitem{b11} R. Shu, X. Gu, and W. Enck, ``A Study of Security Vulnerabilities on Docker Hub,'' in \textit{Proc. ACM CODASPY}, 2017, pp. 269--280.
\bibitem{b12} X. Lin, L. Lei, Y. Wang, J. Jing, K. Sun, and Q. Zhou, ``A Measurement Study on Linux Container Security: Attacks and Countermeasures,'' in \textit{Proc. ACM ACSAC}, 2018, pp. 418--429.
\bibitem{b13} E. G. Young, P. Zhu, T. Caraza-Harter, A. C. Arpaci-Dusseau, and R. H. Arpaci-Dusseau, ``The True Cost of Containing: A gVisor Case Study,'' in \textit{Proc. USENIX HotCloud}, 2019.
\bibitem{b14} A. Randazzo and I. Ferretti, ``Kata Containers: An Emerging Architecture for Enabling MEC Services in Fast and Secure Way,'' in \textit{Proc. IEEE ICITST}, 2019, pp. 1--6.
\bibitem{b15} Snyk Security Research, ``Leaky Vessels: Docker and runc Container Breakout Vulnerabilities (CVE-2024-21626),'' Snyk Blog, Jan. 2024. [Online]. Available: \url{https://snyk.io/blog/leaky-vessels-docker-runc-container-breakout-vulnerabilities/}
\bibitem{b16} M. Myrbakken and R. Colomo-Palacios, ``DevSecOps: A Multivocal Literature Review,'' in \textit{Proc. SWQD}, Springer LNBIP, vol. 302, 2017, pp. 17--29.
\bibitem{b17} A. Shostack, \textit{Threat Modeling: Designing for Security}. Wiley, 2014.
\end{thebibliography}

\end{document}
